\section{Conclusions}\label{sec:conclusions}

This study tests whether post-policy air-quality observations in a local three-district Madrid panel remain below a pre-policy predictive baseline under robustness checks. The central idea is simple: train forecasting models on the pre-policy regime and examine whether the observed post-policy path falls below the no-change forecast after broader model comparison, temporal sensitivity analysis, and placebo checks.

The practical value of the design lies in its discipline. The case study remains focused on three districts and on the combination of air-quality, district traffic intensity, and meteorological information, but the empirical logic is strengthened through Ridge, LightGBM, and LSTM comparisons, naive seasonal benchmarks, windows of 6 and 24 steps, a supplementary 12-step rolling-origin sensitivity, feature ablations, fixed-horizon placebos, and sensitivity to outlier filtering. The results select LightGBM with a 6-step window as the best main-split pre-policy specification by MAE, while the rolling-origin check keeps LightGBM as the preferred family even when the exact short window is not unique. The stacked multivariate LSTM remains weaker than LightGBM overall, which suggests that, in this setting, a strong tabular baseline is more reliable than a more expressive sequential alternative.

Under that preferred specification, the post-policy signal remains positive on average, with a mean PIP of 70.18\% and a positive mean residual of 4.32. The same model also outperforms persistence and seasonal-naive forecasting, its aggregate post-policy summaries remain positive under day-block bootstrap intervals, and its first 60 post-policy days are more favorable than the nearby late-2021 placebo boundaries. At the same time, the new comparison shows why stronger controls matter: Ridge can produce much larger apparent effects, but it does so with substantially worse pre-policy fit. The ablations also indicate that meteorological information matters more than district traffic intensity for preserving the signal in this dataset, while the IQR-based preprocessing mainly affects magnitude rather than direction.

The final claim should therefore remain cautious. The evidence supports a local post-policy deviation that is compatible with lower observed post-policy concentrations relative to the no-change forecast; it does not provide definitive proof of a causal effect, a traffic mechanism, or city-wide generalization. Nearby placebos are weaker than the January 2022 boundary, but they are not null, so some temporal drift remains. The placebo contrast is supportive but not dispositive. That narrower claim is scientifically stronger because it is anchored in explicit model comparison, visible placebo checks, genuine temporal sensitivity, and compact robustness checks rather than in a single favorable architecture.

The code used to generate the empirical results is available at \url{https://github.com/Semllo/Robust-Counterfactual-Forecasting}, together with instructions for retrieving the official Madrid Open Data sources named in Section~\ref{sec:data}.

\subsection*{Funding}
This research received no specific grant from any funding agency in the public, commercial, or not-for-profit sectors.
